% =========================================================
% IV. LEGALSCM-RAG
% =========================================================
\section{LegalSCM-RAG}
\label{sec:legalscm_rag}

We propose \textit{LegalSCM-RAG}, a retrieval-augmented generation framework
for multi-hop legal question answering that integrates retrieval over a
directed condition--consequence graph, multi-hop path reasoning, and selective
verification through a Legal Structural Causal Model (\textit{LegalSCM}).

Given an input question $q$, the overall pipeline is
\begin{equation}
\begin{aligned}
q
&\rightarrow
\text{Causal Retrieval}
\rightarrow
\mathcal{P}(q)
\rightarrow
P^*\\
&\rightarrow
\text{Question-Aware Verification}
\rightarrow
\mathcal{E}^*
\rightarrow
(a,C),
\end{aligned}
\label{eq:overall_pipeline}
\end{equation}
where $P^*$ is the primary reasoning path, $\mathcal{E}^*$ is the selected
evidence, $a$ is the answer, and $C$ is the set of cited provisions.

\subsection{Legal Graph-Based Retrieval}

The system embeds the question and retrieves semantically relevant legal events
and rules from causal memory. Retrieved events serve as seeds for graph
expansion. Starting from each seed, the system traverses directed
condition--consequence relations to construct
\begin{equation}
\mathcal{P}(q)=\{P_1,P_2,\ldots,P_K\},
\end{equation}
where
\begin{equation}
P_i=v_0\xrightarrow{r_1}v_1\xrightarrow{r_2}\cdots
\xrightarrow{r_h}v_h.
\end{equation}
Unlike conventional semantic RAG, the retrieval unit includes legal events and
the directed rule relations connecting them, rather than only independent text
passages.

\subsection{Reasoning Path Ranking}

Candidate paths are ranked using semantic and structural signals:
\begin{equation}
S(P_i,q)=
\alpha S_{\mathrm{sem}}+
\beta S_{\mathrm{graph}}+
\gamma S_{\mathrm{seed}}+
\delta S_{\mathrm{path}}.
\end{equation}
The primary path is
\begin{equation}
P^*=\arg\max_{P_i\in\mathcal{P}(q)}S(P_i,q).
\end{equation}
Rules and provisions associated with $P^*$ are passed to verification as
structured evidence.

\subsection{LegalSCM Fixed-Point Inference}

LegalSCM is a deterministic three-valued normative rule model. Each event is in
$\{\mathrm{TRUE},\mathrm{FALSE},\mathrm{UNKNOWN}\}$. A rule has the form
\begin{equation}
X_1\land\cdots\land X_k\rightarrow Y.
\end{equation}
It fires only when all condition literals have their required settled states
and no encoded exception is active. Conditions within one rule are therefore
conjunctive. Multiple independently activated rules deriving the same state of
$Y$ are alternative, disjunctive support. Opposite supports are retained as a
conflict and mapped to an effective \textsc{unknown} state. The rule operator
is iterated until a fixed point is reached, as formalized in
Section~\ref{sec:problem_formulation}.

The canonical engine supports multiple conditions and exceptions, but the
legacy records used in the reported experiment instantiate one positive
condition and one positive effect per rule. Consequently, the present results
do not constitute an empirical evaluation of multi-condition conjunction or
explicit exception handling.

LegalSCM models normative dependencies extracted from statutes. It neither
discovers causal relations from observational data nor estimates statistical
causal effects.

\subsection{Hard Intervention and Topological Diagnostic}

For an explicit query about removing an intermediate event $M$, LegalSCM
performs a hard intervention relative to the encoded rule model
\cite{pearl2009causality}:
\begin{equation}
do(X_M=0),\qquad 0\equiv\mathrm{FALSE}.
\end{equation}
The intervention fixes $X_M$ to false, blocks every rule mechanism whose effect
is $X_M$, and recomputes the fixed point. If the factual target $Y$ is true and
\begin{equation}
Y_{do(X_M=0)}=\mathrm{FALSE},
\end{equation}
then $M$ is necessary under the encoded context. If
$Y_{do(X_M=0)}=\mathrm{TRUE}$, an alternative mechanism supports the target.
An unknown state in either world yields an indeterminate result.

This operation is distinct from the auxiliary topological ablation that forms
$G^{-M}=G[V\setminus\{M\}]$ and searches for a bounded alternative path. The
latter removes the node and its incident edges and is retained only as a
diagnostic or fallback. It cannot be interpreted as a Pearl-style intervention.
Moreover, LegalSCM performs no exogenous-state abduction; the result is
model-relative necessity in the supplied normative context, not unit-level
actual causation.

\subsection{Question-Aware Verification}

The verifier selects one of three routes from the query. A direct claim
$A\rightarrow C$ is checked against conservatively grounded graph endpoints.
A standard legal question receives factual LegalSCM validation of the primary
path without an intervention. Only an explicit mediator-removal or necessity
query requests $do(X_M=0)$ and compares factual and intervened worlds. This
routing prevents a topological direct-edge test or ordinary factual validation
from being mislabeled as a structural intervention.

\subsection{Selective Verification}

The system distinguishes event states
$\{\mathrm{TRUE},\mathrm{FALSE},\mathrm{UNKNOWN}\}$, internal intervention
statuses such as necessary or non-necessary, and final query labels
\begin{equation}
\mathcal{Z}=\{\textsc{supported},\textsc{reject-direct-claim},
\textsc{uncertain}\}.
\end{equation}
A verifier label becomes authoritative only when confidence, endpoint or
mediator grounding, primary-path coverage, and state-consistency checks pass:
\begin{equation}
z=
\begin{cases}
z_{\mathrm{SCM}}, & G_{\mathrm{commit}}=1,\\[3pt]
z_{\mathrm{LLM}}, & G_{\mathrm{commit}}=0.
\end{cases}
\end{equation}
When the gate fails, the verifier does not force an \textsc{uncertain} label.
Instead, the unchanged Causal Path RAG evidence is supplied to the common
generation model. Full LegalSCM is therefore a selective hybrid rather than an
SCM-authoritative classifier for every sample.

\subsection{Evidence Refinement and Answer Generation}

After verification, the system constructs
\begin{equation}
\mathcal{E}^*=\{r_j\mid r_j\text{ supports }P^*\}
\end{equation}
and provides
\begin{equation}
I_q=(q,P^*,z,\mathcal{E}^*)
\end{equation}
to the language model, which generates
\begin{equation}
(a,C)=f_{\theta}(I_q).
\end{equation}
This separation makes retrieval, rule inference, the selected verification
route, and generation independently auditable.
